\section{Discussion}
\textbf{Scaling Dependence on Evaluation Environment and Metrics.} 
Reinforcement learning optimizes directly for environment rewards~\citep{10.5555/3312046}, which in principle allows unbounded capability—as demonstrated by AlphaZero mastering board games~\citep{silver2017masteringchessshogiselfplay}, AlphaFold predicting protein structures~\citep{jumper2021highly}, and frontier LLMs such as Gemini-2.5-Pro achieving IMO-level performance~\citep{huang2025gemini25procapable}. In contrast, text-based LLMs lack a well-defined RL environment, forcing us to rely on human-curated datasets as proxies. Test loss thus serves as a pragmatic but imperfect metric: it is monotonic and convergent, yet heavily dependent on dataset construction and task difficulty, with different benchmarks (e.g., GSM8K vs. AIME, Section~\ref{subsec:domain-transfer}) showing distinct convergence rates. This task dependence makes the absolute coefficients of our fitted scaling laws ($K_{max}, N_0, E$) difficult to interpret universally. Prior work proposed ``intrinsic performance’’—the minimum compute needed to reach a target reward—as a normalization across environments~\citep{hilton2023scaling}, but we did not find an analogous measure in large-scale LLMs. Establishing principled, environment-independent evaluation protocols remains an open and critical challenge for RL-based scaling studies.

\textbf{Scaling Dependence on Model Scale.} 
While larger models (0.5B-72B) exhibit superior efficiency, our analytic term $k(N)$ confirms this advantage is not infinite. Instead, efficiency gains follow a saturation curve toward a theoretical limit $K_{max}$ (Eq.~\ref{eq:efficiency_intro}), implying diminishing marginal returns at extreme scales.
This finding implies that scaling up models beyond a certain point, while still yielding absolute performance gains, suffers from diminishing marginal returns in efficiency.
% Our study examined models from 0.5B to 14B parameters, revealing consistent gains in sample and compute efficiency for larger models under RL post-training. This range, however, is still below frontier deployments, and our larger-scale experiments are underway. We anticipate that the observed advantages hold when tasks are of moderate difficulty and data are well curated, but may not extend indefinitely. On simple benchmarks such as GSM8K, performance tends to saturate as models grow, while for extremely hard problems (e.g., NP-hard tasks or open research challenges) neither small nor large models exhibit meaningful convergence. A plausible hypothesis is the existence of a critical model scale: large enough to encode broad commonsense knowledge for efficient RL adaptation, but not so over-parameterized that returns diminish. Identifying such a “sweet spot’’ remains an important direction for future work.


\textbf{Dependence on RL Algorithm.}
Our analysis is based on GRPO, a mainstream and stable RL post-training algorithm for LLMs that uses an actor-only design and normalizes rewards across responses. Comparative study with alternative RL algorithms ~\citep{cui2025entropymechanismreinforcementlearning} reports minor differences in training curves. Whether more advanced algorithms can significantly improve sample efficiency or stability—and thereby reshape the scaling frontier—remains an important open question.

\textbf{Future of LLM Agent.}
The integration of reinforcement learning with agentic LLMs is increasingly viewed as a promising direction~\citep{zhang2025landscapeagenticreinforcementlearning, zhang2025surveyreinforcementlearninglarge}. Both theoretical and empirical studies show that augmentations such as external tool use and long-term memory can substantially boost model performance~\citep{lin2025understandingtoolintegratedreasoning,houliston2025provablebenefitsintoollearning,mai2025agentrlscalinglaw}. In the multi-agent setting, reward-driven self-organization~\citep{zhou2025reso} and co-evolutionary interaction rewards~\citep{xue2026comascoevolvingmultiagentsystems} further show that collaborative LLM agents can surpass individual model performance. We anticipate that such agentic mechanisms will markedly improve the scaling behavior of RL-trained LLMs: by offloading deterministic computations to tools and focusing learning on high-level decision making, these models could achieve much higher efficiency, effectively shifting the performance frontier upward for a given compute or data budget. Understanding the scaling laws of these agentic systems is, therefore, a key and exciting avenue for future research.
